\newpage
\chapter{Техники за одговор на колизија}

Детекцијата на колизија од претходната глава дава 3 информации на фазата за одговор:
\begin{itemize}
	\item \textbf{Длабочина на пенетрација} \textit{d} - колку длабоко навлегла честичката во телото
	\item \textbf{Нормала на точката на контакт} $\vec{n}$ - нормализирана насока
	\item \textbf{Точка на контакт} (најчесто иста со моменталната позиција на честичката)
\end{itemize}

Во фазата на одговор треба:
\begin{itemize}
	\item Да се одделат телата, односно да доведе \textit{d} на 0 со поместување на едно од телата
	\item Да се промени нормалната брзина, односно да сведе на нула (нееластично), да се сврти обратно (еластично), или некаде помеѓу со коефициентот на еластичност
	\item Да се промени тангентната брзина, односно да се намали со триење со што се намалува кинетичката енергија
\end{itemize}

Има четири парадигми за решавање на проблемот кои се разликуваат во одлуката кога и како се прават промените

\subsection{Методи со пенали}
Најстар и најинтуитивен пристап\cite{hahn1988realistic}, се третира продирањето како компресирана пружина и се применува сила на реставрирање $F = -k \cdot d \cdot \vec{n}$, со термин на намалување $F_d = -c \cdot(\vec{v} \cdot \vec{n})\vec{n}$ за да се намали нормалната брзина. Силата влегува во равенките за движење како сила на телото кај контактот, интегрирана со временскиот чекор што го користи системот.

Тие се лесни за имплементација, може да се користат во било кој тип на динамика (експлицитен Ојлер, RK4, полуимплицитен итн.) и се диференцијабилни, што е важно за апликации базирани на градиенти. Но имаат проблем со цврстината, за да се намали продирањето, \textit{k} мора да биде големо, но големо \textit{k} ја дестабилизира интеграцијата, па на наидува на проблем на видлива мекост за мало \textit{k} и експлодирање на симулацијата за големо \textit{k}. Друг проблем е што може да се внесе енергија во системот и што има осцилација околу точката на контакт поради користењето на пружини.

Овие методи се важни историски но не се користат во денешните апликации.

\subsection{Импулсни методи}

Наместо да се примени сила врз конечен временски интервал, се применува моментен импулс \textit{J} на контактната точка што директно ја менува брзината. Големината на импулсот се пресметува од :
$$(\vec{v}_i' - \vec{v}_j') \cdot \vec{n} = -e \cdot (\vec{v}_i - \vec{v}_j) \cdot \vec{n}$$

\noindent каде $\vec{v}_{i}^{'} = \vec{v}_i + J \cdot \frac{\vec{n}}{m_i}$ и $e$ е коефициентот на реституција.

Предностите им се што немаат параметар за цврстина и реституцијата е едноставна физичка мерка од [0,1] и за еден контакт има само една пресметка, без интегрирање. Тие се стандард во енџини за цврсти тела. Слабостите им се што кога има повеќе истовремени контакти има лошо скалирање на цената за пресметка. Тие се одлични за ретки контактни точки кај цврсти тела, но неа за густи и повеќебројни контакти кај меките тела

\subsection{Контакт базиран на ограничувања}
Се третира секој контакт како ограничување со нееднаквост $d(\vec{x}) \geq 0$ и да се реши ограничувањето заедно со сите други ограничувања во системот. Триењето влегува како дополнително ограничување на тангентната брзина.



\subsection{Проекција на позиции (геометриска корекција во решавањето)}
Откако ќе се решат еластичните ограничувања, се поминува низ секоја честичка и оние кои што имаат навлезено во колизиско тело се проектираат назад на површината, а тангенцијалниот дел од движењето се анулира. Потоа брзината се пресметува од делтата на позиции, што автоматски го поништува движењето навнатре и го рефлектира триењето.

Ова е наједноставен метод што добро функционира, потребни се само неколку линии код. Методот е безусловно стабилен и лесно се интегрира со PBD/XPBD без менување на решавачот. 